\documentclass[12pt,a4paper]{ctexart}

% 版面
% 对齐 Word 的 A4 默认版式：左右 3.17cm，上下 2.54cm
\usepackage[left=3.17cm,right=3.17cm,top=2.54cm,bottom=2.54cm]{geometry}
\usepackage{amsmath,amssymb}
\usepackage{booktabs}
\usepackage{graphicx}
\usepackage{float}
\usepackage{caption}
\usepackage{subcaption}
\usepackage{xcolor}
\usepackage[hidelinks]{hyperref}
\usepackage{enumitem}
\usepackage{siunitx}
\usepackage[numbers]{gbt7714}

\linespread{1.5}\selectfont
\captionsetup{font=small,labelsep=quad}
\setlist[itemize]{leftmargin=1.4em,itemsep=2pt,topsep=3pt}
\setlist[enumerate]{leftmargin=1.6em,itemsep=2pt,topsep=3pt}

\newcommand{\todo}[1]{\textcolor{red}{[待补：#1]}}
% 引用编号整体不可断行，也禁止在编号之前断行（否则 [4] 会掉到下一行行首）
\let\origcite\cite
\renewcommand{\cite}[2][]{\nobreak\mbox{\origcite[#1]{#2}}}

\title{\textbf{面向小型无人机场景的动作条件隐空间预测与局部路径修正}}
\author{
  刘子瑄 \quad 马婧怡 \quad 盛祺 \quad 张子睿 \\[6pt]
  {\normalsize 华南理工大学}
}
\date{}

\begin{document}
\setcitestyle{numbers,sort&compress}
\bibpunct{[}{]}{,}{n}{,}{,}
\pagestyle{plain}
\thispagestyle{empty}

% 标题区：手写排版，不用 \maketitle，避免过大留白
\begin{center}
  {\LARGE\bfseries 面向小型无人机场景的动作条件隐空间预测与局部路径修正\par}
  \vspace{16pt}
  {\large 刘子瑄\quad 马婧怡\quad 盛祺\quad 张子睿\par}
  \vspace{6pt}
  {\small\kaishu 华南理工大学\par}
\end{center}
\vspace{10pt}

% 摘要
\begingroup\small
\noindent\textbf{摘要：}小型无人机在地块分散、视野受限的场景中飞行时，需要在看清障碍之前判断``接下来哪些区域会被占据''，才能及时调整路径。现有视频预测研究多以被动观测为设定，只预测画面如何变化，不把飞控动作作为输入，因而无法回答``如果执行某个动作会怎样''。

本文以仿真为数据来源：在 Unreal Engine 中构建采集管线，同步记录相机画面、相机位姿、飞控动作与占据真值；训练一个可训练参数 4.7\,M 的隐空间预测模型\cite{assran2023ijepa,bardes2024vjepa}，以 8 帧观测与候选动作作为输入，输出未来 1 秒的占据分布，并以指数滑动平均目标编码器与方差/协方差正则抑制表征退化。模型被接入分层规划：全局路径由 A* 依据静态地图生成，模型仅修正在途的局部路径，使世界模型的贡献与全局规划解耦。

在同一全局路径下对照两种安全判据，以模型预测代替当前帧判断后，碰撞次数由 \todo{X} 降至 \todo{Y}（下降 \todo{Z\%}）。正则强度的三点对照进一步表明其与下游精度呈非单调关系；当正则过弱时，隐空间预测误差仍接近最优，而下游指标全部为零，说明该类误差不能单独作为评估依据。

\vspace{6pt}
\noindent\textbf{关键词：}隐空间预测；动作条件；无人机避障；分层规划
\endgroup

\vspace{16pt}


% 引言
\section{引言}

无人机在小型农业与园区巡检中的应用逐渐增多。此类场景地块面积小、分布分散、边界不规则，大型农机难以适配，作业主要依赖小型多旋翼飞行器。在这类环境中，飞行器需要在视野受限的条件下判断``接下来哪些区域会被占据''，以支持避障与路径调整。

现有的视频预测方法多针对被动观测，回答的是``接下来画面会如何变化''\cite{bardes2024vjepa,hafner2024dreamerv3}，而不回答``如果执行某个动作会如何变化''。后者才是决策所需的信息，其建模要求将动作作为显式输入。

将动作引入预测模型，存在三个具体困难。

\textbf{其一，缺少动作条件的数据。}公开的无人机视觉数据集以检测与跟踪为目标构建\cite{zhu2021visdrone}，只提供图像与目标框，不含相机位姿与飞控指令，无法支撑动作条件预测的训练与评估。

\textbf{其二，缺少有判别力的评估指标。}在持续性主导的场景中，全图 IoU 会被大量静止区域稀释。本文在 VisDrone 上的实测显示，``未来等于现在''的平凡基线在预测 5 步后仍取得 0.406 的全图 IoU，该指标不足以判别模型是否学到动力学。

\textbf{其三，小模型下的表征坍缩。}隐空间预测存在退化解：当编码器对所有输入输出同一向量时，预测误差为零而表征不含信息。抑制该解需要两个约束，即目标编码器的指数滑动平均（EMA）与表征的方差/协方差正则（VICReg）\cite{bardes2022vicreg}。

针对上述问题，本文以仿真为数据来源构建动作条件世界模型，并将其接入分层规划框架。工作内容如下：

\begin{enumerate}
  \item \textbf{数据}：在 Unreal Engine 中建立采集管线，同步记录相机画面、位姿、飞控动作与占据真值，形成具有因果结构的数据集；
  \item \textbf{模型}：训练可训练参数 4.7\,M 的隐空间预测模型，输入为 8 帧观测与候选动作，输出未来 1 秒的占据分布，以 EMA 目标编码器与 VICReg 抑制坍缩；
  \item \textbf{系统}：全局路径由 A* 依据静态地图生成，世界模型仅修正在途的局部路径，使``世界模型的价值''与``全局规划的价值''解耦；
  \item \textbf{评估}：在同一全局路径下，对照``以当前帧判断安全性''与``以模型预测判断安全性''的碰撞次数。
\end{enumerate}

本文全部实验在仿真中完成，不包含真机验证，也不包含风扰、通信延迟等真实扰动。

% 相关工作
\section{相关工作}

\subsection{隐空间预测}
以表征为预测目标的思路源于联合嵌入预测架构（JEPA）\cite{assran2023ijepa}：编码器将观测映射为表征，预测器在表征空间中推断被遮挡部分的目标，训练目标是最小化预测表征与目标表征之间的距离。V-JEPA 将该框架扩展到视频\cite{bardes2024vjepa}，以指数滑动平均（EMA）维护的目标编码器提供训练目标，梯度不回传至目标分支。

不在像素空间预测的理由有两方面。其一，像素值受纹理与光照主导，误差被大面积静态区域稀释；其二，重建高保真画面所需的解码能力与预测动力学并不等价。但表征预测存在退化解：当编码器对所有输入输出同一向量时，预测误差为零而表征不含信息。VICReg\cite{bardes2022vicreg}以方差项与协方差项对该退化解加以约束，前者要求每一维的标准差不低于阈值，后者惩罚维度之间的线性相关。

\subsection{无人机视觉导航中的世界模型}
Dreamer 系列以循环状态空间模型学习隐空间动力学，并在隐空间中进行想象以训练策略\cite{hafner2024dreamerv3,hafner2020dreamer}。面向四旋翼平台，MAD\cite{liu2026mad}指出以原始图像重建为自监督目标并非最优，转而重建以机体为中心的占据图与可见性图以及本体状态，使隐状态直接编码局部几何、可见性历史与自运动，并在真实森林环境中以 5.05\,m/s 完成实机飞行。AirDreamer\cite{liu2026airdreamer}在世界模型提供的环境表征之上训练强化学习策略，以不含人工塑形的稀疏奖励避免局部最优并鼓励偏航控制，在困难地图上的导航成功率较最优基线高 5.3\%，且从仿真到实机部署无需微调。

另一条路线将世界模型的作用限制在训练阶段。FlowPilot\cite{flowpilot2026}联合去噪未来深度表征与轨迹控制点，部署时不生成未来观测而直接输出轨迹；AeroAct\cite{aeroact2026}以未来帧作为密集监督训练动作头，推理阶段完全跳过视频生成。二者共同的理由是生成式预测的延迟与误差累积不适合进入飞行的实时回路。

上述工作多以强化学习或模仿学习得到策略。本文不训练策略，而是保留显式的预测模块，由规划器在预测结果上做决策，因而需要单独度量``预测是否比当前观测更有用''。

\subsection{学习模型与规划}
将学习到的动力学模型嵌入模型预测控制（MPC）是成熟的做法\cite{nagabandi2018neural,chua2018pets}：在每一控制周期以当前状态为初值，对未来若干步的候选动作序列做预测并择优，仅执行首步。该结构与本文采用的``全局路径 + 局部修正''一致，区别在于本文的预测目标为占据分布而非低维状态，且预测模型由自监督方式训练，不依赖专家轨迹。

\subsection{仿真数据}
无人机视觉研究普遍依赖仿真器采集带标注的数据\cite{shah2018airsim}。可编程渲染引擎能够同步输出画面、位姿与几何真值，这一点对动作条件预测尤为关键：公开的真实航拍数据集通常只提供图像与目标框，缺少与画面逐帧对齐的控制指令。

% 方法
\section{方法}

\subsection{数据采集}
\todo{250 字：场景构成、采集内容（画面/位姿/动作/占据）、动作定义、分辨率与帧率}

\subsection{模型结构}

输入为 8 帧 192$\times$108 的 RGB 画面。每一帧被切分为 $6\times6$ 的图像块并线性投影为 576 个 token\cite{dosovitskiy2021vit}，由 6 层、3 头的 Transformer 编码为 192 维表征。块边长取 6 像素是由目标尺度反推的：真实数据中目标的中位边长为 7.5 像素（以 192 宽为基准），块边长大于该值时一个块内会混入多个目标，表征无法区分相邻个体。

预测器为 4 层时序 Transformer，采用块因果注意力：第 $i$ 帧的 token 只能看到第 $0$ 至 $i$ 帧，因而可以在推理时逐帧自回归展开而不泄漏未来信息；其输出为下一帧的 576 个 token。

教师网络为编码器副本，参数按 $\tau=0.996$ 的指数滑动平均更新，不接收梯度。探针为三层卷积（$1\times1$、$3\times3$、$1\times1$），将 token 网格双线性上采样为 $27\times48$ 的占据图；训练探针时编码器冻结。

模型可训练参数为 4.70\,M（编码器 2.80\,M、预测器 1.87\,M、探针 0.03\,M），另有 2.80\,M 的教师参数不参与优化，合计 7.50\,M。单卡 8\,GB 显存下取批大小 3、梯度累积 3。

\subsection{训练目标}
预测损失定义为预测表征与目标表征之间的余弦距离：
\begin{equation}
\mathcal{L}_{\text{pred}} = \frac{1}{K}\sum_{k=1}^{K}\Bigl(1-\cos\bigl(\hat{z}_{t+k},\,z_{t+k}\bigr)\Bigr),
\end{equation}
并以 VICReg\cite{bardes2022vicreg} 的方差项与协方差项约束表征：
\begin{equation}
\mathcal{L} = \mathcal{L}_{\text{pred}} + \lambda_{\text{var}}\mathcal{L}_{\text{var}} + \lambda_{\text{cov}}\mathcal{L}_{\text{cov}}.
\end{equation}
预测项在每一步预测上计算余弦距离，教师分支不回传梯度。方差项要求每一维的标准差不低于 1，用以阻止表征退化为常数；协方差项惩罚维度之间的线性相关，用以减少维度冗余。二者权重取 $\lambda_{\text{var}}=0.10$、$\lambda_{\text{cov}}=0.01$，其中 $\lambda_{\text{var}}$ 由三点对照确定，依据见第 4.2 节。

\subsection{规划接口}
给定在上帝视角下选定的起点与终点，A* 在静态地图上生成全局路径。该地图只包含场景中的固定结构，不含移动目标，因而全局路径不随目标运动而改变。飞行过程中每 0.5\,s 触发一次局部修正：以最近 8 帧画面与当前期望动作序列为输入，模型输出未来 5 步的占据分布；若预测占据与全局路径的横向间距低于安全阈值，则对路径施加横向偏移并重新采样参考点。修正只改变路径的局部形状，不改变终点，因而全局规划的最优性与世界模型的贡献可以分开度量。

% 实验
\section{实验}

\subsection{设置}
\todo{150 字：数据规模、训练配置、评估指标定义}

\subsection{正则强度的影响}
在真实数据上对正则权重 $\lambda_{\text{var}}$ 做三点对照，结果见表~\ref{tab:lambda}。权重过小时表征坍缩，隐空间余弦接近 1 而下游指标全部归零；权重过大时表征过于分散，下游精度下降。中间取值在综合指标上最优。

\begin{table}[H]
\centering
\caption{正则权重对照（真实数据，预测时域 5 步）。$\lambda_{\text{var}}=0.05$ 时探针不产生任何正预测，下游指标为零。空白格表示该基线不存在对应定义。}
\label{tab:lambda}
\begin{tabular}{lccccc}
\toprule
$\lambda_{\text{var}}$ & 闭环余弦 & $R_{\text{new}}$ & $P_{\text{occ}}$ & F1 & IoU \\
\midrule
0.05 & 0.905 & 0.000 & 0.000 & 0.000 & 0.000 \\
0.10 & 0.640 & 0.349 & 0.412 & \textbf{0.378} & \textbf{0.190} \\
0.15 & 0.306 & \textbf{0.507} & 0.179 & 0.265 & 0.127 \\
\midrule
平凡基线 &              & 0.000 &              &              & 0.406 \\
\bottomrule
\end{tabular}
\end{table}

\subsection{主结果}
\todo{250 字 + 表：碰撞次数对照（用当前帧 vs 用模型预测）}

\subsection{消融}
\todo{150 字：有动作 / 无动作 / 恒等基线三组}

% 讨论
\section{讨论}

\subsection{隐空间误差不能作为评价指标}
与文献\cite{liu2026mad}的观察一致，以几何或占据为目标优于以像素为目标。表~\ref{tab:lambda} 中，$\lambda_{\text{var}}=0.05$ 的模型取得最高的闭环余弦（0.905），但其下游指标全部为零。这说明隐空间预测误差在表征坍缩时会给出与真实能力相反的排序，评估必须以下游任务为准。

\subsection{增益来源}
\todo{200 字：动作条件带来的增益来自何处；与``用当前帧''的差异分析}

% 局限
\section{局限与展望}

本文的工作有以下边界：全部实验在仿真中完成，未经真机验证；动作不含噪声，未建模通信延迟与风扰；预测时域为 1 秒，未涉及多机协同与长时域规划；仅建模飞行动作，未涉及农事作业动作。

后续工作的一个方向是多模态融合。实际系统中观测将来自多光谱、热红外与地面传感器，其采样率与坐标系并不一致。隐空间预测不要求像素级对齐，为该类融合提供了接口。

% 结论
\section{结论}

\todo{100 字：三条结论}

% 参考文献
\begingroup\small
\bibliographystyle{gbt7714-numerical}
\bibliography{refs}
\endgroup

\end{document}
